\documentclass[11pt]{article}

\newcommand{\cnum}{Math 245B}
\newcommand{\ced}{Winter 2026}
\newcommand{\ctitle}[4]{\title{\vspace{-0.5in}\cnum, \ced\\Assignment #1: #2}\author{\vspace{-0.35in}\\#3, #4}}
\usepackage{enumitem}
\usepackage[usenames,dvipsnames,svgnames,table,hyperref]{xcolor}
\usepackage{amsmath, amsfonts}
\usepackage{graphicx} % Required for inserting images
\usepackage{float}
\usepackage{listings}
\usepackage{xcolor}
\usepackage{hyperref}
\usepackage{comment}

\renewcommand*{\theenumi}{\alph{enumi}}
\renewcommand*\labelenumi{(\theenumi)}
\renewcommand*{\theenumii}{\roman{enumii}}
\renewcommand*\labelenumii{\theenumii.}

\DeclareMathOperator{\spt}{spt}

\definecolor{codegreen}{rgb}{0,0.6,0}
\definecolor{codegray}{rgb}{0.5,0.5,0.5}
\definecolor{codepurple}{rgb}{0.58,0,0.82}
\definecolor{backcolour}{rgb}{0.95,0.95,0.92}
\definecolor{header}{HTML}{205459}
\definecolor{solution}{HTML}{204d52}

\newcommand{\solution}[1]{{{\textcolor{header}{Solution:} \textcolor{solution}{#1}}}}
\setlist[enumerate]{label=(\alph*)}

\lstdefinestyle{mystyle}{
    backgroundcolor=\color{backcolour},
    commentstyle=\color{codegreen},
    keywordstyle=\color{magenta},
    numberstyle=\tiny\color{codegray},
    stringstyle=\color{codepurple},
    basicstyle=\ttfamily\footnotesize,
    breakatwhitespace=false,
    breaklines=true,
    captionpos=b,
    keepspaces=true,
    numbers=left,
    numbersep=5pt,
    showspaces=false,
    showstringspaces=false,
    showtabs=false,
    tabsize=2
}


\begin{document}
\ctitle{\#2}{}{Asher Christian}{006-150-286}
\date{}
\maketitle

\section{Exercise 2.1}
\emph{
    Let $\mu $ be a Borel (outer) measure on a separable metric space $X$.
    Let $V$ be the union of all open sets $O \subset X$ such that $\mu(O) = 0$
    (spt $\mu  := X \setminus V$ is called the support of $\mu $ ).
    \begin{enumerate}
        \item Show that $X$ is second countable (has a countable basis of open sets)
        \item Show that $\mu(X \setminus \spt \mu) = 0$
        \item Further assume that $\mu(X) < \infty$. Show that $\spt \mu$ is the smallest closed subset $K$ such that
             $\mu(K) = \mu(X)$.
    \end{enumerate}
}
\solution{
    \begin{enumerate}
        \item
            Since $X$ is separable there exists a countable dense subset $Q \subset X$ such that for all $x \in X, \epsilon > 0, \exists q \in Q, d(q,x) < \epsilon$.
            Let $Q = \{q_i\}_{i=1}^{\infty}$ and Define
            \[
                B := \{D(q_i,\frac{1}{m}) : i,m \in \mathbb{N}\}
            \] 
            where $D(x,r)$ is the open ball in  $X$ centered at  $x$ of radius $r$. I claim that $B$ is a base of $X$.
            Indeed let  $O \subset X$ be open then for each $x \in O$ there exists  $\epsilon(x) >0$ such that $D(x,\epsilon(x)) \subset O$,
            and there exists $q(x) \in Q$ such that  $d(x,q(x)) < \frac{\epsilon(x)}{3}$ and $m(x)$ such that  $\frac{\epsilon(x)}{3} \le \frac{1}{m(x)} < \frac{\epsilon}{2}$ together we get that 
            \[
                O = \bigcup_{x \in O}D(q(x),\frac{1}{m(x)})
            \] 
            since each $x \in O$ is in the ball  $D(q(x),\frac{1}{m(x)})$ which itself is contained in $D(x,\epsilon(x)) \subset O$. Also $B$ is a countable collection of sets since  $Q$ is countable, so $B$ is a countable base. 
        \item 
            we have
            \[
            X \setminus \spt \mu = X \setminus (X \setminus V) = V
            \] 
            \[
                V = \bigcup_{\substack{O \subset X \\ \mu(O) = 0}}O = \bigcup_{\substack{O \subset X\\ \mu(O) = 0}}\bigcup_{\substack{B_i \in B\\ B_i \subset O}}B_i = \bigcup_{j \in M \subset \mathbb{N}}^{}B_j
            \] 
            for some collection $M$ of elements of the basis from the last question $B$. Then we have by countable subadditivity of measure
            \[
                \mu(V) = \mu(\bigcup_{j \in M}B_j) \le \sum_{j \in M}^{}\mu(B_j) = 0
            \] 
            since each $B_j \subset O$ so $\mu(B_j) \le \mu(O) = 0$.
        \item Assume $\mu(X) < \infty$ then we have
            \[
            \mu(\spt \mu) = \mu(X \setminus V) = \mu(X) - \mu(V) = \mu(X)
            \] 
            Assume for the sake of contradiction that there exists $K \subset \spt \mu$ $K \ne \spt \mu$ obeying $\mu(X) = \mu(K)$
            then
            \[
            \mu(X \setminus K) = \mu(X) - \mu(K) = 0
            \] 
            so $K^{c}$ is an open set with measure 0 and $V \subset K^{c}$, however $K^{c}$ is open and so $K^{c} \subset V$ thus $K^{c} = V$ and $K = V^{c} = \spt \mu$ thus $\spt \mu$ is the smallest closed subset that has full measure.
    \end{enumerate}
}

\section{Exercise 2.4}
\emph{
    Let $\mathcal{L}$ be the class of Lebesgue measurable sets in  $ \mathbb{R}^{d}$ and let $\mu$
    be the Lebesgue measure on $ \mathbb{R}^{d}$. Let $(\nu_j)_j$ be a monotone sequence of Radon measures on  $ \mathbb{R}^{d}$ such
    that $\nu(E) := \lim_{j\to \infty}\nu_j(E)$ is a Radon measure on $ \mathbb{R}^{d}$. Show that
    \[
    D_\mu\nu = \lim_{j\to \infty} D_\mu\nu_j \quad a.e.
    \] 
}
\solution{
    Assume first that the sequence $(\nu_j)_j$ is monotone increasing, that is for all  $A \subset \mathbb{R}^{d}$ $\nu_j[A] \le \nu_{j+1}[A]$ then we have similarly $\nu_j[A] \le \nu[A]$ and so for all $r$
\[
    \frac{\nu(B_r(x))}{\mu(B_r(x))} \ge \frac{\nu_j(B_r(x))}{\mu(B_r(x))} \implies D_\mu \nu \ge D_\mu \nu_j
\] 
taking limits we get
\[
D_\mu \nu \ge \lim_{j\to \infty} D_\mu \nu_J
\] 
We know via the lebesgue decomposition theorem that there exists a set $B \subset \mathbb{R}^{d}$ such that $\mu[B^{c}] = 0$ and
\[
    \nu = \nu|_B + \nu|_{B^{c}} = \nu_{abs} + \nu_s \qquad \nu_{abs} << \mu \qquad \nu_{s} \perp \mu
\] 
There also exists a sequence of sets $(B_j)_j \subset 2^{\mathbb{R}^{d}}$ such that $\mu[B_j^{c}] = 0$ for all $j$
\[
    \nu_j = \nu_j|_{B_j} + \nu_j|_{B_J^{c}}
\] 
with the absolutely continuity and singularity holding as well. Then $\mu[(B \cap \bigcap_{j=1}^{\infty}B_j)^{c}] = 0$
so it suffices to show the inequality on $\tilde B = B \cap \bigcap_{j=1}^{\infty}B_j$ since if it holds almost everywhere on  $\tilde B$ then it holds $\mu$-almost everywhere. Let $E \subset \tilde B$ be a $\mu$-measurable set then
\[
    \int_{E}^{} D_\mu \nu(x) d \mu = \nu[E] = \lim_{j\to\infty }\nu_j[E] = \lim_{j\to\infty } \int_{E}^{}D_\mu \nu_j(x) d \mu = \int_{E}^{} \lim_{j\to \infty} D_\mu \nu_j(x) d \mu
\] 
Firstly by lebesgue differentiation theorem since when we restrict $\nu_j$ and  $\nu$ to  $\tilde B$ they are absolutely continuous. Further the other inequality is ipmlied by the monotone convergence theorem since $D_\mu \nu_j$ is a monotone increasing sequence of functions with respect to $j$
\[
\int_{E}^{}D_\mu \nu d \mu = \int_{E}^{} D_\mu \nu_j d \mu
\] 
for all $E \subset \tilde B$ which implies that
\[
D_\mu \nu =  \lim_{j\to \infty} D_\mu \nu_j 
\] 
$\mu$-almost everywhere on $\tilde B$ which further implies that  $D_\mu \nu = D_\mu \nu_j$ $\mu$-almost everywhere since $\tilde B^{c}$ is a null set.
For the other possiblility, the $(\nu_j)_j$ are monotone decreasing, keeping $\tilde B$ defined as before, we further restrict ourselves to proving that for all  $E \subset \tilde B$
 \[
\int_{E \cap B_n(0)}^{}D_\mu \nu(x)d \mu = \int_{E \cap B_n(0)}^{} \lim_{j\to \infty} D_\mu \nu_j(x) d \mu
\] 
after which, we conclude that the two functions agree on all bounded sets and so the set of which they agree almost sure is the union of $E \cap B_n(0)$ over all  $n \in \mathbb{N}$ and we ultimately conclude that they agree almost everywhere on all of $ \mathbb{R}^{d}$.
To prove this equality we cite the dominated convergence theorem to state that for fixed $E$
 \[
\int_{E \cap B_n(0)}^{}D_\mu \nu_j(x) d \mu \le \int_{ E \cap B_n(0)}^{} D_\mu \nu_1(x) d \mu < \infty
\] 
With the inequality holding pointwise as well.
since $\nu$ is a Radon measure so the integrals are finite on all compact sets. Thus by dominated convergence we argue
\[
    \int_{E \cap B_n(0)}^{} D_\mu \nu(x) d \mu = \lim_{j\to \infty}\int_{E \cap B_n(0)}^{}D_\mu \nu_j(x) d \mu = \int_{E \cap B_n(0)}^{}\lim_{j\to \infty}D_\mu \nu_j d \mu
\] 
So we conclude as before that the functions are the same almost everywhere.
}

\section{Exercise 2.6}
\emph{
    Let $a < b$ be real numbers and let $F : [a,b] \rightarrow \mathbb{R}$ be a continuous functin
    \begin{enumerate}
        \item Show that the set $A$ of points of differentiablility of $F$ is a Borel set.
        \item Suppose that $F$ is monotone nondecreasing and $\mu_F$ is the associated Lebesgue-Stieltjes measure.
            Let $B_{\infty} := \{F' = +\infty\}$. Show that for any Borel set  $E$ 
            \[
                \mu_F(E) = \int_{E}^{}F'(t)dt + \mu_F(E \cap B_{\infty})
            \] 
    \end{enumerate}
    Similar results can be obtained when $F$ is of bounded variations and $\mu_F$ is a signed measure
}

\end{document}
